%% revision_literatura_bistatica.tex
%% Revisión de literatura: SAR biestático (criterio imprescindible de selección)
%%               + búsqueda dirigida: optimización basada en ángulo de Brewster
%%               para maximizar potencia recibida (biestático y monoestático)
%% Uso previsto: insumo para el capítulo de estado del arte / trabajos relacionados
%%               de la tesis doctoral (doc/dissertacao_tomo_sar.tex)
%%
%% Compilar: pdflatex revision_literatura_bistatica.tex (dos pasadas)

\documentclass[11pt,a4paper]{article}

% ─── paquetes ────────────────────────────────────────────────────────────────
\usepackage[utf8]{inputenc}
\usepackage[T1]{fontenc}
\usepackage[spanish]{babel}
\usepackage{lmodern}
\usepackage{geometry}
\geometry{left=2.7cm, right=2.5cm, top=2.7cm, bottom=2.7cm}

\usepackage{amsmath,amssymb}
\usepackage{booktabs}
\usepackage{array}
\usepackage{longtable}
\usepackage{enumitem}
\usepackage{xcolor}
\usepackage[colorlinks=true,linkcolor=blue!50!black,citecolor=blue!50!black,urlcolor=blue!50!black]{hyperref}
\usepackage{titlesec}

\titleformat{\section}{\normalfont\Large\bfseries}{\thesection}{1em}{}
\titleformat{\subsection}{\normalfont\large\bfseries}{\thesubsection}{1em}{}

\setlength{\parindent}{0pt}
\setlength{\parskip}{0.5\baselineskip}

\title{\bfseries Revisión de literatura: sistemas SAR biestáticos\\
\large Estudios afines al presente trabajo doctoral\\
(criterio imprescindible de selección: configuración biestática)}
\author{David Atencia}
\date{Septiembre de 2026}

\begin{document}
\maketitle

\section{Introducción y criterio de selección}

Este documento reúne 22 trabajos de la literatura que abordan, en distintos
niveles de la cadena de procesado SAR (geometría, resolución, formación de
imagen, tomografía 3D, sistemas experimentales y explotación del ángulo de
Brewster), el caso en que el transmisor y el receptor \textbf{no comparten
la misma plataforma} —condición \emph{imprescindible} de selección para los
bloques A--D, ya que es la que define el marco de trabajo de la presente
tesis: SAR biestático helicoidal con propagación en uno o varios medios
(aire--suelo). El bloque E amplía la búsqueda al caso monoestático de forma
deliberada, por tratarse de una comprobación de novedad específica (véase
\S\ref{sec:brewster}).

Se priorizaron publicaciones recientes (2018--2025) para capturar el estado
del arte actual, complementadas con tres referencias fundacionales
(Arikan~\&~Munson 1989, Ding~\&~Munson 2002, Horne~\&~Yates 2002) que ya
forman parte del marco teórico de esta tesis y de la carpeta
\texttt{papers/} del repositorio, incluidas aquí para contextualizar la
evolución del campo. Cada entrada indica una síntesis del abstract original
y el beneficio concreto de leerla/implementarla para el trabajo desarrollado
en \texttt{simulations/} y \texttt{doc/dissertacao\_tomo\_sar.tex}.

Los trabajos se agrupan en cinco bloques temáticos:
\begin{enumerate}[itemsep=2pt]
  \item Fundamentos geométricos y resolución espacial biestática (\S\ref{sec:fundamentos}).
  \item Formación de imagen y algoritmos de reconstrucción biestática (\S\ref{sec:algoritmos}).
  \item Tomografía SAR biestática y multiestática 3D (\S\ref{sec:tomografia}).
  \item Sistemas biestáticos aerotransportados y UAV (\S\ref{sec:uav}).
  \item Ángulo de Brewster y optimización de potencia recibida, biestático y monoestático (\S\ref{sec:brewster}).
\end{enumerate}

% ══════════════════════════════════════════════════════════════════════════
\section{Fundamentos geométricos y resolución espacial biestática}
\label{sec:fundamentos}

\begin{enumerate}[label=\textbf{[A\arabic*]}, leftmargin=2.2cm, itemsep=10pt]

\item Horne, A.M.; Yates, G. (2002). \textit{Bistatic synthetic aperture
radar.} En \emph{Proc. RADAR 2002}, Edimburgo, IEE Publication No.\,490,
pp.\,6--10. \cite{horne2002}\\
\textbf{Síntesis del abstract:} el SAR es cada vez más importante en
vigilancia terrestre militar por su capacidad de operar en cualquier clima,
de día y de noche, y de detectar, clasificar y geolocalizar objetos a largo
alcance. El biestático (Tx/Rx en plataformas separadas) se ve como un medio
de contrarrestar la vulnerabilidad electrónica, manteniendo la plataforma
receptora encubierta mientras el transmisor —más vulnerable— puede
desplegarse más lejos. El artículo repasa la investigación en curso de
QinetiQ sobre las técnicas de procesado para SAR biestático, los problemas
fundamentales que introduce (sincronización, ruido de fase, ejes de imagen
no ortogonales) y cómo superarlos.\\
\textbf{Beneficio:} da el marco conceptual completo que sostiene el modelo
geométrico biestático de esta tesis; leerlo ayuda a anticipar problemas
prácticos (sincronización Tx/Rx, ruido de fase acumulado) antes de un
eventual despliegue real con dos UAV.

\item Homer, J.; Donskoi, E.; Mojarrabi, B.; Palmer, J.; Kubik, K. (2005).
\textit{Three-dimensional bistatic synthetic aperture radar imaging system:
spatial resolution analysis.} \emph{IEE Proceedings -- Radar, Sonar and
Navigation}, 152(6), 391--394. \cite{homer2005}\\
\textbf{Síntesis del abstract:} presentan un marco matemático para analizar
la capacidad de resolución espacial de un sistema SAR biestático 3D
general, en el que transmisor, receptor y objeto (o región de interés)
pueden ser no coplanares y en el que el transmisor y/o el receptor pueden
estar en movimiento. Con aplicaciones de observación de la Tierra en mente,
el marco se usa para caracterizar cuantitativamente la capacidad de
resolución 2D en el plano del terreno del sistema de imagen 3D; los
resultados teóricos se confirman mediante simulación.\\
\textbf{Beneficio:} es la referencia teórica más cercana a la derivación de
$\delta_{xy}$, $\delta_z$ de esta tesis pero para geometría arbitraria (no
solo circular/espiral); permite generalizar el análisis de resolución a
trayectorias de vuelo no ideales (deriva, viento) y validar
cuantitativamente la fórmula de resolución en el plano de terreno con un
caso 3D genérico.

\item Garza, G.; Qiao, Z. (2010). \textit{Resolution analysis of bistatic
SAR.} En \emph{Proc. SPIE 8021, Radar Sensor Technology XV}, 80210F.
\cite{garza2010}\\
\textbf{Síntesis del abstract:} analizan la resolución de imágenes BISAR de
objetos estacionarios reconstruidas mediante un método de inversión por
backprojection filtrado, derivado de un modelo de ecuación de onda escalar,
que permite tener en cuenta los efectos de los patrones de haz de la antena
y de trayectorias de vuelo arbitrarias. El análisis se realiza examinando
la variedad (\emph{manifold}) de recolección de datos para distintas
geometrías experimentales y parámetros del sistema.\\
\textbf{Beneficio:} su función de dispersión puntual (PSF), ligada
explícitamente al patrón de haz de la antena, es aplicable para refinar el
modelo de resolución de esta tesis incorporando el diagrama de radiación
real de la antena Tx/Rx (hoy en \texttt{getAntennaGain.m}), que las
fórmulas actuales asumen idealizado.

\item Chen, Y.; Chen, R.; Liu, H.; Guo, J.; Wang, Y.; Zhang, J. (2022).
\textit{Research on configuration constraints of airborne bistatic SARs.}
\emph{Sensors}, 22(17), 6534.
\href{https://doi.org/10.3390/s22176534}{doi:10.3390/s22176534}
\cite{chen2022}\\
\textbf{Síntesis del abstract:} a partir del análisis de la geometría de
imagen del SAR biestático aerotransportado, emplean un algoritmo de chirp
scaling no lineal extendido para simular y verificar el efecto de imagen.
Proponen un modelo de resolución 2D biestático basado en teoría de
gradientes y analizan las restricciones de los parámetros de trayectoria de
vuelo multiplataforma que satisfacen la precisión de imagen requerida.
Mediante simulación de imagen biestática con trayectorias de vuelo
cooperativas en distintos escenarios, revelan la envolvente de restricción
de configuración espacial entre los vehículos para lograr la resolución
óptima.\\
\textbf{Beneficio:} su ``envolvente de configuración válida'' entre Tx y Rx
es una herramienta directamente aplicable para acotar el espacio de
búsqueda del optimizador \texttt{run\_plano\_de\_voo.m}/\texttt{\_multislab.m}
(hoy basado en multi-start sin restricción geométrica derivada
analíticamente), lo que podría acelerar su convergencia.

\end{enumerate}

% ══════════════════════════════════════════════════════════════════════════
\section{Formación de imagen y algoritmos de reconstrucción biestática}
\label{sec:algoritmos}

\begin{enumerate}[label=\textbf{[B\arabic*]}, leftmargin=2.2cm, itemsep=10pt]

\item Arikan, O.; Munson, D.C.\,Jr. (1989). \textit{A tomographic
formulation of bistatic synthetic aperture radar.} En Porter, W.A.; Kak,
S.C. (eds.), \emph{Advances in Communications and Signal Processing},
Lecture Notes in Control and Information Sciences, vol.\,129, pp.\,418--431,
Springer. (Presentado en \emph{Proc. ComCon '88}, Baton Rouge, LA, oct.\,1988.)
\cite{arikan1989}\\
\textbf{Síntesis del abstract:} un radar biestático es aquel con antenas
transmisora y receptora separadas por una distancia comparable a la
distancia al objetivo; el concepto no es nuevo pero fue eclipsado durante
décadas por el monoestático. En años recientes ha habido interés militar
considerable en el SAR biestático de modo spotlight (BSSAR), donde un
transmisor de alta potencia opera a distancia segura mientras un receptor
encubierto vuela cerca del área objetivo. Dan una derivación básica de
BSSAR desde primeros principios, mostrando —de forma análoga al caso
monoestático (MSSAR)— que puede explicarse mediante el teorema de
corte-proyección de la tomografía asistida por computador (CAT), sin
depender del efecto Doppler: a cada ángulo de mira, los datos BSSAR
demodulados son muestras aproximadamente banda-limitadas de una rebanada de
la transformada de Fourier 2D de la reflectividad del objetivo. Examinan la
ubicación de las muestras en el dominio de Fourier y la curvatura del
frente de onda para varios casos especiales de movimiento de Tx y Rx.\\
\textbf{Beneficio:} es la base tomográfica (dominio-$k$) que sostiene el
Cap.~2--3 de la disertación; releerlo con foco en la curvatura del frente
de onda ayuda a verificar que la aproximación de campo lejano usada en
\texttt{precomputeTxDataMultislab.m} sigue siendo válida para las
distancias Tx--blanco del escenario UAV de esta tesis.

\item Ding, Y.; Munson, D.C.\,Jr. (2002). \textit{A fast back-projection
algorithm for bistatic SAR imaging.} En \emph{Proc. IEEE Int. Conf. on Image
Processing (ICIP)}, vol.\,II, pp.\,449--452. \cite{ding2002}\\
\textbf{Síntesis del abstract:} usando un modelo de campo lejano, el SAR
biestático adquiere datos de Fourier en una grilla no cartesiana poco
habitual. Los algoritmos previos de formación de imagen se basaban
principalmente en reconstrucción directa de Fourier (DFR) para aprovechar
la FFT, pero la cobertura irregular de los datos disponibles y la
interpolación 2D en el dominio de Fourier pueden afectar la precisión de la
reconstrucción. El backprojection (BP) evita estos problemas pero tiene un
costo computacional enorme, $O(N^3)$ frente a $O(N^2\log_2 N)$ de la DFR.
Presentan un algoritmo BP rápido para SAR biestático, motivado por un
algoritmo rápido de BP propuesto recientemente para tomografía, que logra
un costo computacional del mismo orden que la DFR y que además puede
usarse para imagen de campo cercano; validan el rendimiento mediante
simulación.\\
\textbf{Beneficio:} es el antecedente algorítmico directo del backprojection
implementado en \texttt{simulations/proc/}; portar su versión rápida
(jerárquica) permitiría escalar el pipeline PROC a grillas de procesamiento
volumétricas más finas sin el costo $O(N^3)$ actual, beneficio directo de
rendimiento para \texttt{run\_snell\_volumetrico\_pipeline.m}.

\item Zhang, H.; Tang, J.; Wang, R.; Deng, Y.; Wang, W.; Li, N. (2018).
\textit{An accelerated backprojection algorithm for monostatic and bistatic
SAR processing.} \emph{Remote Sensing}, 10(1), 140.
\href{https://doi.org/10.3390/rs10010140}{doi:10.3390/rs10010140}
\cite{zhang2018}\\
\textbf{Síntesis del abstract:} proponen un algoritmo BP acelerado a partir
de Block\_FFBP en el que un número fijo de ``pivotes'' —en vez de centros de
haz que crecen con cada etapa de factorización— construye la relación de
retardo de propagación, y los datos de rango se procesan como un bloque
completo sin partición en la dimensión de rango, gracias a un esquema
integrado de procesamiento de rango. El esquema se extiende a SAR
biestático (BiBulk\_FFBP) y monoestático (MoBulk\_FFBP), transfiriendo el
cálculo de rango oblicuo del sistema polar al cartesiano para simplificarlo.
Validan el algoritmo con datos reales del experimento SAR biestático
espacial/estacionario TerraSAR-X en modo staring spotlight, y con un
experimento SAR spotlight aerotransportado de 2016.\\
\textbf{Beneficio:} ruta de aceleración computacional validada con datos
reales de una misión biestática espacial; referencia de eficiencia
aplicable a la escalabilidad del backprojection volumétrico de
\texttt{run\_snell\_volumetrico\_pipeline.m} si se busca portarlo a un
esquema jerárquico (FFBP).

\item Wang, Z.; Liu, M.; Ai, G.; Wang, P.; Lv, K. (2020). \textit{Focusing of
bistatic SAR with curved trajectory based on extended azimuth nonlinear
chirp scaling.} \emph{IEEE Transactions on Geoscience and Remote Sensing},
58(6), 4160--4179. \cite{wang2020}\\
\textbf{Síntesis del abstract:} proponen un modelo de rango de movimiento
de alto orden y un algoritmo de chirp scaling no lineal acimutal extendido
(EANLCS) con corrección residual del centroide Doppler, para enfocar datos
SAR biestáticos con trayectoria curva de las plataformas. El trabajo aborda
el acoplamiento rango-acimut adicional y la variación espacial de la
migración de celda de rango y de la tasa de modulación Doppler causadas por
la trayectoria curva, efectos que los algoritmos tradicionales —derivados
para trayectorias rectilíneas de velocidad uniforme— no pueden procesar.\\
\textbf{Beneficio:} es la referencia más cercana al caso de trayectorias
helicoidales cónicas de esta tesis, ya que ninguno de los algoritmos
``clásicos'' biestáticos (Arikan, Ding) asume trayectoria curva; revisar su
modelo de rango de alto orden es un candidato de mejora de precisión de
fase para el procesador si en el futuro se simplifica el cálculo exacto de
distancia vía Fermat que usa hoy el pipeline.

\item Sun, H.; Sun, Z.; Chen, T.; Miao, Y.; Wu, J.; Yang, J. (2022).
\textit{An efficient backprojection algorithm based on wavenumber-domain
spectral splicing for monostatic and bistatic SAR configurations.}
\emph{Remote Sensing}, 14(8), 1885.
\href{https://doi.org/10.3390/rs14081885}{doi:10.3390/rs14081885}
\cite{sun2022}\\
\textbf{Síntesis del abstract:} en el FBP tradicional, cada subapertura usa
su propio sistema de coordenadas polares centrado en el centro de dicha
subapertura, por lo que las imágenes de subapertura deben proyectarse a un
sistema de coordenadas uniforme antes de poder sumarse coherentemente,
demandando mucho cálculo. Proponen usar el mismo sistema de coordenadas
polares para todas las subaperturas y empalmar (\emph{splice}) sus imágenes
en el dominio del número de onda, sumándolas directamente tras
sobremuestreo, evitando la proyección a un sistema uniforme y mejorando así
precisión y eficiencia; el método es aplicable a configuraciones mono y
biestáticas.\\
\textbf{Beneficio:} técnica de fusión en dominio-$k$ directamente
relacionada con el análisis de $k_z$/$k_{xy}$ de
\texttt{notes/modelo\_conceptual.md}; útil para acelerar el post-procesado
de imagen volumétrica sin sacrificar la precisión de fase que determina la
resolución $\delta_z$ de esta tesis.

\item Zhao, T.; Wang, J.; Cheng, Z.; Huang, Z.; Song, J. (2025). \textit{A
subimage autofocus bistatic ground Cartesian back-projection algorithm for
passive bistatic SAR based on GEO satellites.} \emph{Remote Sensing}, 17(9),
1576. \href{https://doi.org/10.3390/rs17091576}{doi:10.3390/rs17091576}
\cite{zhao2025}\\
\textbf{Síntesis del abstract:} el SAR biestático pasivo (PBSAR) basado en
señales de satélites geoestacionarios (GEO) muestra potencial significativo
para imagen de alta resolución, pero enfrenta desafíos duales de
eficiencia computacional y compensación de errores de fase. Proponen un
algoritmo de backprojection Cartesiano en tierra biestático con
autoenfoque de subimagen (SIAF-BGCBP), mediante corrección del espectro de
número de onda de subimagen, que reduce en un 90\,\% la complejidad de la
estimación de fase al optimizar la densidad de píxeles de subimagen,
manteniendo la precisión de estimación.\\
\textbf{Beneficio:} relevante si en el futuro se explora una arquitectura
biestática pasiva de bajo costo (iluminador fijo/satelital + receptor UAV)
como alternativa al esquema Tx/Rx activo dual actual; su autoenfoque por
subimagen es transferible al problema de corrección de fase por errores de
posicionamiento GPS del receptor durante el vuelo.

\end{enumerate}

% ══════════════════════════════════════════════════════════════════════════
\section{Tomografía SAR biestática y multiestática 3D}
\label{sec:tomografia}

\begin{enumerate}[label=\textbf{[C\arabic*]}, leftmargin=2.2cm, itemsep=10pt]

\item Ge, N.; Zhu, X.X. (2019). \textit{Bistatic-like differential SAR
tomography.} \emph{IEEE Transactions on Geoscience and Remote Sensing},
57(8), 5883--5893. \cite{ge2019}\\
\textbf{Síntesis del abstract:} abordan la tomografía SAR diferencial
(D-TomoSAR) en áreas urbanas usando adquisiciones espaciales biestáticas o
``pursuit monostatic''. Los interferogramas biestáticos/pursuit-monostatic
no sufren decorrelación temporal significativa ni pantalla de fase
atmosférica, por lo que son ideales para la reconstrucción de elevación.
Proponen un marco de dos etapas: primero reconstruyen la elevación de los
dispersores con interferogramas biestáticos-\emph{like}, y luego incorporan
esa elevación como prior determinístico para estimar la deformación a
partir de interferogramas convencionales de paso repetido
(\emph{repeat-pass}).\\
\textbf{Beneficio:} su estrategia de ``dos etapas'' (elevación biestática
como prior + refinamiento con datos repeat-pass) es transferible para
separar el problema de localización 3D del blanco enterrado (elevación o
profundidad) del problema de estimación de reflectividad, potencialmente
reduciendo el espacio de búsqueda del backprojection volumétrico de esta
tesis.

\item Chen, Z.; Li, Y.; Hu, C. et al. (2024). \textit{Repeat-pass
space-surface bistatic SAR tomography: accurate imaging and first
experiment.} \emph{Science China Information Sciences}, 67, 192304.
\href{https://doi.org/10.1007/s11432-024-4089-2}{doi:10.1007/s11432-024-4089-2}
\cite{chenz2024}\\
\textbf{Síntesis del abstract:} la tomografía SAR biestática
espacio-superficie (SS-BiSAR) ofrece un ángulo de observación adicional al
SAR monoestático espacial, siendo prometedora para la recuperación de
deformación de alta precisión en regiones locales. Presentan el primer
experimento de tomografía SS-BiSAR de paso repetido, usando el satélite
gemelo chino LuTan-1 en banda L, adquiriendo 11 imágenes de pista repetida
y recuperando nubes de puntos 3D de un área edificada con
$\sim$4\,m de precisión. Proponen un método de corrección de efemérides
basado en la transformada chirp-Z para corregir la distorsión geométrica de
la imagen y los errores de fase interferométrica causados por
imprecisiones orbitales.\\
\textbf{Beneficio:} primera demostración experimental de tomografía
biestática 3D real en banda L —misma banda que la configuración
\texttt{\_EMIRADOS} de este proyecto—; su método de corrección de
``efemérides''/posición de plataforma es directamente análogo al problema
de errores de posicionamiento GPS del receptor UAV que afectaría la fase
del backprojection de esta tesis en campo.

\item Wu, J.; Yang, J.; Li, Z.; Pu, W.; Sun, Z.; An, H.; Yang, Q. (2025).
\textit{Review of bistatic synthetic aperture radar imaging methods.}
\emph{Journal of Radars}, 14(5), 1115--1141.
\href{https://doi.org/10.12000/JR25067}{doi:10.12000/JR25067}
\cite{wu2025}\\
\textbf{Síntesis del abstract:} el sistema SAR biestático emplea
plataformas de transmisión y recepción espacialmente separadas para
proveer imágenes de alta resolución de escenas y blancos terrestres y
marítimos en entornos complejos, con ventajas de configuración flexible,
fuerte capacidad de ocultamiento, alta resistencia a interferencia y
adquisición más completa de información del blanco. El procesado de imagen
es crítico para obtener imágenes BiSAR de alta resolución, ya que el
modelo de eco y las características del BiSAR difieren sustancialmente de
las del SAR monoestático tradicional. La revisión examina desafíos y
soluciones clave para varias configuraciones BiSAR —aerotransportado, de
plataformas de alta velocidad y maniobrabilidad, espacial heterogéneo y
espacial homogéneo—, junto con compensación de movimiento e imagen de
blancos móviles, repasa avances de investigación nacionales e
internacionales, y da una perspectiva de tendencias futuras.\\
\textbf{Beneficio:} es el mapa más actualizado (2025) del campo completo
del SAR biestático; leerlo íntegramente es la forma más eficiente de
verificar que el capítulo de estado del arte de la disertación no omite
ninguna línea de investigación activa (p.\,ej. BiSAR espacial
homogéneo/heterogéneo) antes de la defensa.

\item He, J.; Xie, H.; Liu, H.; Wu, Z.; Xu, B.; Zhu, N.; Lu, Z.; Qin, P.
(2025). \textit{Multi-baseline bistatic SAR three-dimensional imaging
method based on phase error calibration combining PGA and EB-ISOA.}
\emph{Remote Sensing}, 17(3), 363.
\href{https://doi.org/10.3390/rs17030363}{doi:10.3390/rs17030363}
\cite{he2025}\\
\textbf{Síntesis del abstract:} la tomografía SAR (TomoSAR) multi-línea de
base monoestática sufre decorrelación temporal en escenas como bosques;
además, el \emph{jitter} de plataforma y los errores de medición de
posición contaminan los datos multi-línea de base con errores de fase que
degradan la calidad de la reconstrucción. Proponen un método de imagen 3D
biestático multi-línea de base que combina la calibración de fase PGA
(\emph{Phase Gradient Autofocus}) con un nuevo método EB-ISOA para corregir
estos errores en un escenario biestático.\\
\textbf{Beneficio:} aborda exactamente el problema de ``jitter de
plataforma + error de posición degradando la fase'' que enfrentaría
cualquier implementación real de esta tesis con dos UAV volando
trayectorias helicoidales independientes; su combinación PGA+EB-ISOA es un
candidato concreto de autoenfoque a incorporar en el pipeline PROC antes
de una validación experimental en campo.

\item Andre, D.; Krishnan, V.P.; Nolan, C. (2025). \textit{Artefact analysis
of multistatic three-dimensional SAR imaging with two linear
trajectories.} arXiv:2503.11323. \cite{andre2025}\\
\textbf{Síntesis del abstract:} con la mejora de la tecnología de radar, el
SAR multiestático es hoy una posibilidad realista, por ejemplo en
constelaciones satelitales o sistemas de aeronaves no tripuladas, lo que
plantea la necesidad de investigar geometrías multiestáticas útiles. Dan un
análisis microlocal de una geometría de adquisición de datos multiestática
para imagen SAR 3D, en la que el transmisor y el receptor se mueven
independientemente a lo largo de dos líneas rectas que cubren una región de
interés (ROI). Usando técnicas microlocales, muestran cómo la geometría de
adquisición determina la presencia de artefactos en la imagen resultante y
cómo evitarlos; su aporte principal es una condición de ``time-gating''
sobre los datos que garantiza la ausencia de artefactos en la ROI, y la
verifican con varias simulaciones numéricas bajo una formulación
independiente del tiempo.\\
\textbf{Beneficio:} herramienta analítica rigurosa para predecir y evitar
artefactos de imagen antes de simular; aplicar su condición de time-gating
al pipeline GS$\to$PROC de esta tesis permitiría filtrar a priori qué
combinaciones de posición Tx/Rx en la trayectoria helicoidal aportan
energía ``limpia'' al blanco, reduciendo artefactos en el backprojection
volumétrico sin necesidad de barridos de prueba y error.

\end{enumerate}

% ══════════════════════════════════════════════════════════════════════════
\section{Sistemas biestáticos aerotransportados y UAV}
\label{sec:uav}

\begin{enumerate}[label=\textbf{[D\arabic*]}, leftmargin=2.2cm, itemsep=10pt]

\item Liu, M. et al. (2024). \textit{Unmanned airborne bistatic
interferometric synthetic aperture radar data processing method using
bi-directional synchronization chain signals.} \emph{Remote Sensing},
16(5), 769. \href{https://doi.org/10.3390/rs16050769}{doi:10.3390/rs16050769}
\cite{liu2024}\\
\textbf{Síntesis del abstract:} los UAV equipados con InSAR biestático
tienen costo relativamente bajo y alta flexibilidad, útiles para mapeo y
exploración de recursos terrestres. Debido a desafíos como la
sincronización espaciotemporal y los errores de movimiento, proponen un
método integral de procesamiento de datos para InSAR biestático UAV de
pequeño porte, integrando señales de cadena de sincronización
bidireccional, compensación de errores de sincronización temporal y de
fase, refinamiento de trayectoria, imagen InSAR biestático de alta
precisión y procesamiento interferométrico. Los resultados muestran que el
método mejora la precisión de medición interferométrica en puntos de
calibración de 0.66\,m a 0.42\,m.\\
\textbf{Beneficio:} es un sistema UAV biestático real y completo (hardware
+ procesado) con el mismo desafío de sincronización de dos plataformas
aéreas independientes que enfrentaría una implementación física de esta
tesis; su técnica de sincronización bidireccional es un candidato directo
de diseño de hardware para una futura validación experimental del
escenario Tx/Rx helicoidal.

\item Grathwohl, A.; Meinecke, B.; Widmann, M.; Kanz, J.; Waldschmidt, C.
(2023). \textit{UAV-based bistatic SAR-imaging using a stationary
repeater.} \emph{IEEE Journal of Microwaves}, 3(2), 625--634.
\href{https://doi.org/10.1109/JMW.2023.3253667}{doi:10.1109/JMW.2023.3253667}
\cite{grathwohl2023}\\
\textbf{Síntesis del abstract:} el SAR es una técnica establecida de imagen
de radar de alta resolución que se apoya en el automovimiento del radar y
en un largo tiempo de integración coherente, iluminando el blanco desde
distintos ángulos de observación; usar una configuración biestática para
adquirir simultáneamente múltiples ángulos de observación incrementa aún
más estos beneficios. Proponen un sistema SAR biestático usando un
repetidor activo estacionario (Bistatic Repeater-SAR, BR-SAR), donde el
repetidor actúa como segundo observador proporcionando una vía biestática
con bajo costo y esfuerzo de hardware; describen las consideraciones de
procesado de señal y hardware necesarias, validando el enfoque con
imágenes monoestáticas, biestáticas y multiestáticas combinadas de alta
resolución.\\
\textbf{Beneficio:} demuestra que un sistema biestático UAV puede
construirse con hardware simple (repetidor pasivo/estacionario) en vez de
dos plataformas activas sincronizadas —alternativa de bajo costo a
considerar para una futura validación experimental de campo de la
trayectoria de receptor optimizada por \texttt{run\_plano\_de\_voo.m}.

\item Kanz, J.; Gesell, C.; Bonfert, C.; Werbunat, D.; Grathwohl, A.;
Aguilar, J.; Vossiek, M.; Waldschmidt, C. (2025). \textit{UAV-borne digital
radar system for coherent multistatic SAR imaging.} arXiv:2507.20792.
\cite{kanz2025}\\
\textbf{Síntesis del abstract:} los avances en conversores
analógico-digitales (ADC) permiten adoptar arquitecturas de radar digital
que muestrean directamente la señal de radiofrecuencia, eliminando la
necesidad de downconversión analógica y dando mayor flexibilidad en el
diseño de forma de onda, en particular mediante esquemas de modulación
digital como OFDM. Presentan un sistema de radar digital montado en un UAV
que emplea formas de onda OFDM para imagen SAR multiestática coherente en
banda L: un nodo primario transmite y adquiere datos monoestáticos,
mientras nodos secundarios operan solo en modo recepción, capturando tanto
la señal reflejada por la escena como una señal de enlace lateral directo.
El sistema mantiene coherencia tanto en tiempo rápido como en tiempo lento,
esencial para la imagen multiestática; al ser pasivos, los nodos
secundarios permiten escalar fácilmente a un enjambre mayor de UAV.
Detallan el flujo de procesamiento completo, incluyendo el análisis y
corrección de errores de sincronización derivados de la operación
desacoplada de los nodos, y validan el método con mediciones estáticas y
un experimento SAR biestático UAV que produce imágenes monoestáticas,
biestáticas y multiestáticas combinadas de alta resolución.\\
\textbf{Beneficio:} arquitectura de radar digital OFDM con nodo Tx activo +
nodos Rx pasivos escalables es la ruta más directa para extender el
escenario ``Tx fijo + Rx optimizado'' de esta tesis a validación de campo
real; su solución de sincronización mediante enlace lateral directo
(\emph{sidelink}) resuelve el mismo problema de sincronización de fase
Tx/Rx que sería crítico para trasladar
\texttt{run\_plano\_de\_voo\_multislab.m} de simulación a vuelo real.

\end{enumerate}

% ══════════════════════════════════════════════════════════════════════════
\section{Ángulo de Brewster y optimización de potencia recibida (biestático y monoestático)}
\label{sec:brewster}

Esta sección responde a una búsqueda dirigida y específica: \emph{¿existe en
la literatura algún trabajo de SAR tomográfico que optimice la trayectoria
o geometría de adquisición en función del ángulo de Brewster para maximizar
la potencia recibida, ya sea en configuración biestática o monoestática?}
Deliberadamente se relaja aquí el criterio ``imprescindiblemente biestático''
de las secciones A--D para cubrir también el caso monoestático, tal como se
solicitó.

\textbf{Resultado de la búsqueda:} no se encontró ningún trabajo que combine
(i) formación de imagen SAR \emph{tomográfica} (reconstrucción 3D) con (ii)
optimización activa de la trayectoria o geometría Tx/Rx en función del
ángulo de Brewster para maximizar la potencia recibida —ni en biestático ni
en monoestático—. Los trabajos más cercanos encontrados se dividen en dos
categorías, ninguna de las cuales cierra esa brecha:

\begin{enumerate}[label=\textbf{[E\arabic*]}, leftmargin=2.2cm, itemsep=10pt]

\item Ouchi, K.; Yang, C.-S. (2017). \textit{Direct evidence and
implications of Brewster's angle damping observed by synthetic aperture
radar images.} En \emph{Proc. IGARSS 2017.} \cite{ouchi2017}\\
\textbf{Síntesis del abstract:} reportan evidencia directa del
amortiguamiento (\emph{damping}) por ángulo de Brewster y sus implicaciones
en análisis polarimétrico, observado en imágenes SAR de estructuras de
concreto como espigones y rompeolas.\\
\textbf{Categoría:} monoestático, \emph{observación pasiva} del efecto
Brewster en imágenes ya adquiridas (no hay optimización de trayectoria ni
tomografía 3D).\\
\textbf{Beneficio:} confirma que el efecto Brewster es medible y
significativo en SAR monoestático de superficie con datos reales
(TerraSAR-X); útil como evidencia empírica independiente de que el
contraste de potencia por ángulo de incidencia que explota esta tesis en
biestático (Cap.~6) es un fenómeno físicamente observable y no solo
teórico.

\item Ouchi, K.; Yang, C.-S. (2018). \textit{Brewster Angle Damping
Observed in the TerraSAR-X Synthetic Aperture Radar Images of Man-Made
Targets.} \emph{IEEE Geoscience and Remote Sensing Letters}, 15(4),
532--536. \cite{ouchi2018}\\
\textbf{Síntesis del abstract:} versión de revista, ampliada, del trabajo
anterior: analizan el amortiguamiento por ángulo de Brewster observado en
imágenes SAR co-polarizadas de TerraSAR-X sobre blancos artificiales
(estructuras de concreto), cuantificando el fenómeno y discutiendo sus
implicaciones para la interpretación polarimétrica de imágenes SAR de
superficie.\\
\textbf{Categoría:} monoestático, análisis polarimétrico \emph{a
posteriori} (no hay optimización de trayectoria ni reconstrucción 3D).\\
\textbf{Beneficio:} aporta un modelo cuantitativo del amortiguamiento por
Brewster en datos SAR reales, útil como referencia de contraste para
validar experimentalmente —con datos de campo, si estuvieran disponibles—
la ganancia de potencia predicha por el modelo de Fresnel/Brewster usado en
\texttt{objective\_function\_multislab.m}.

\item Thirion-Lefevre, L.; Guinvarc'h, R. (2018). \textit{The double
Brewster angle effect.} \emph{Comptes Rendus Physique}, 19(1--2), 43--53.
\href{https://doi.org/10.1016/j.crhy.2018.02.003}{doi:10.1016/j.crhy.2018.02.003}
\cite{thirion2018}\\
\textbf{Síntesis del abstract:} el efecto de doble ángulo de Brewster (DBE)
extiende el concepto de ángulo de Brewster a la doble reflexión sobre dos
superficies dieléctricas ortogonales, resultante de la combinación de dos
pseudo-ángulos de Brewster que ocurren en dominios de ángulo de incidencia
complementarios. Demuestran teórica y experimentalmente que el DBE puede
observarse potencialmente en cualquier imagen de radar monoestático de la
superficie terrestre.\\
\textbf{Categoría:} monoestático, fenómeno físico de doble reflexión (no
hay optimización de trayectoria ni SAR tomográfico).\\
\textbf{Beneficio:} generaliza el concepto de ángulo de Brewster simple
(usado en esta tesis para la interfaz plana aire-suelo) al caso de
esquinas/superficies dieléctricas ortogonales; relevante como línea de
trabajo futura si se extiende el modelo de esta tesis a geometrías de
subsuelo no planas (zanjas, esquinas de estructuras enterradas).

\item Orfeo, D.; Zhang, Y.; Burns, D.; Miller, J.S.; Huston, D.R.; Xia, T.
(2019). \textit{Bistatic antenna configurations for air-launched ground
penetrating radar.} \emph{Journal of Applied Remote Sensing}, 13(2),
027501. \href{https://doi.org/10.1117/1.JRS.13.027501}{doi:10.1117/1.JRS.13.027501}
\cite{orfeo2019}\\
\textbf{Síntesis del abstract:} el GPR es valioso para detectar objetos
subsuperficiales con poco o ningún contenido metálico (plásticos,
cerámicas, tuberías de concreto), pero los efectos de los parámetros de
configuración de antena —como altura y ángulo— no están bien estudiados
para todas las aplicaciones de sensado. Realizan simulaciones de GPR y
experimentos de laboratorio para evaluar los efectos del ángulo y la
altura de antena en la sensibilidad de un GPR biestático de lanzamiento
aéreo, en la búsqueda de objetos no metálicos enterrados. Los resultados
dan una guía para el desarrollo de sistemas GPR de lanzamiento aéreo
instalados en UAV para escaneo subsuperficial en vuelo.\\
\textbf{Categoría:} \textbf{biestático}, optimización de ángulo/altura de
antena orientada a UAV —el más cercano en espíritu a la optimización de
esta tesis—, pero sin formación de imagen SAR/tomográfica: es un único par
Tx/Rx fijo por medición (``point sensing''), no una síntesis de apertura
con trayectoria continua, y no menciona expresamente el ángulo de Brewster
como criterio de diseño (aunque el efecto físico que exploran —mínimo de
reflexión en la interfaz según ángulo y polarización— es el mismo
fenómeno).\\
\textbf{Beneficio:} es el antecedente más próximo al problema concreto de
esta tesis (GPR biestático UAV con geometría Tx/Rx variable para maximizar
señal transmitida al subsuelo); su validación experimental de laboratorio
es una referencia útil para diseñar un experimento de campo que confirme
la ganancia de potencia predicha por el término $P_r(\theta)$ del costo
$J$ en \texttt{objective\_function\_multislab.m}, y su ausencia de
formación de imagen 3D subraya la brecha que esta tesis cierra.

\end{enumerate}

\textbf{Conclusión de la búsqueda dirigida:} el uso del ángulo de Brewster
como criterio físico para maximizar la transmisión de potencia hacia el
subsuelo está bien establecido en la literatura de GPR de punto único
(antena fija a ángulo/altura óptimos) y se ha \emph{observado} de forma
pasiva en imágenes SAR monoestáticas de superficie. Sin embargo, en
septiembre de 2026 no se encontró ningún trabajo —biestático ni
monoestático— que integre ese criterio dentro de un \emph{optimizador de
trayectoria} acoplado a una \emph{formación de imagen SAR tomográfica 3D}
de blancos subsuperficiales, que es exactamente la contribución de
\texttt{run\_plano\_de\_voo.m}/\texttt{run\_plano\_de\_voo\_multislab.m} en
esta tesis (función de costo
$J=-(1-\alpha)\log_{10}P_r+\alpha\log_{10}(\delta_{xy}\delta_z)$). Esto
refuerza la conclusión de \S\ref{sec:sintesis}: la combinación específica
de SAR biestático tomográfico + refracción multicapa + optimización de
trayectoria por ángulo de Brewster no tiene, hasta donde permite establecer
esta búsqueda, antecedente directo en la literatura.

% ══════════════════════════════════════════════════════════════════════════
\section{Síntesis: posicionamiento de esta tesis respecto a la literatura}
\label{sec:sintesis}

\begin{longtable}{@{}p{3.6cm}p{3.2cm}p{3.2cm}p{5.3cm}@{}}
\toprule
\textbf{Trabajo} & \textbf{Geometría Tx/Rx} & \textbf{Medio de propagación} & \textbf{Diferencia con esta tesis} \\
\midrule
\endhead
Horne \& Yates 2002 / Arikan \& Munson 1989 & Lineal / spotlight & Aire & Sin refracción ni trayectoria helicoidal \\
Homer et al. 2005 / Garza \& Qiao 2010 & General 3D & Aire & Resolución general, sin caso subsuperficial \\
Wang et al. 2020 / Chen et al. 2022 & Curva (aeronave) & Aire & Trayectoria curva pero sin interfaz de dos medios \\
Ge \& Zhu 2019 / Chen Z. et al. 2024 / He et al. 2025 & Multi-baseline / repeat-pass & Aire & Tomografía 3D biestática, pero orientada a superficie/deformación, no a blancos enterrados \\
Liu 2024 / Grathwohl 2023 / Kanz 2025 & UAV, dos plataformas & Aire & Sistemas biestáticos UAV reales, sin modelo de refracción Snell \\
Zhao et al. 2025 & Pasivo (GEO) & Aire & Biestático pasivo, geometría fija del iluminador \\
Wu et al. 2025 (revisión) & Todas & Aire & Mapa de estado del arte, no aborda subsuelo \\
Ouchi \& Yang 2017/2018 / Thirion-Lefevre \& Guinvarc'h 2018 & Monoestático & Aire (observación de superficie) & Observan el efecto Brewster \emph{a posteriori}, sin optimizar trayectoria ni reconstruir en 3D \\
Orfeo et al. 2019 & Biestático, par fijo & Aire + suelo (interfaz plana) & Optimiza ángulo/altura de antena para un punto de medición único, sin síntesis de apertura ni imagen 3D \\
\textbf{Esta tesis} & \textbf{Helicoidal cónica Tx/Rx} & \textbf{Aire + suelo (N capas, Snell/Brewster)} & \textbf{Único trabajo que combina SAR biestático 3D con refracción multicapa y optimización de trayectoria del receptor por ángulo de Brewster} \\
\bottomrule
\end{longtable}

Ningún trabajo revisado combina simultáneamente: (i) geometría biestática
con trayectorias helicoidales cónicas independientes para Tx y Rx, (ii)
propagación con refracción en interfaz aire--suelo (o multicapa), y
(iii) optimización conjunta de potencia recibida (ángulo de Brewster) y
resolución 3D mediante síntesis de apertura SAR/tomográfica. La búsqueda
dirigida de \S\ref{sec:brewster} confirma que ni siquiera relajando la
condición biestática (es decir, incluyendo el caso monoestático) existe
antecedente que cierre esta brecha. Esta combinación constituye la
contribución diferencial de la presente tesis dentro del panorama de la
literatura SAR biestática actual.

% ══════════════════════════════════════════════════════════════════════════
\section*{Referencias}
\addcontentsline{toc}{section}{Referencias}

\begin{thebibliography}{99}

\bibitem{andre2025} D.~Andre, V.~P.~Krishnan, C.~Nolan, ``Artefact Analysis of Multistatic Three-Dimensional SAR Imaging with Two Linear Trajectories,'' arXiv:2503.11323, 2025.

\bibitem{arikan1989} O.~Arikan, D.~C.~Munson~Jr., ``A tomographic formulation of bistatic synthetic aperture radar,'' in \emph{Advances in Communications and Signal Processing}, W.~A.~Porter, S.~C.~Kak, Eds., Lecture Notes in Control and Information Sciences, vol.~129, pp.~418--431, Springer, 1989.

\bibitem{chen2022} Y.~Chen, R.~Chen, H.~Liu, J.~Guo, Y.~Wang, J.~Zhang, ``Research on Configuration Constraints of Airborne Bistatic SARs,'' \emph{Sensors}, vol.~22, no.~17, art.~6534, 2022. doi:10.3390/s22176534

\bibitem{chenz2024} Z.~Chen, Y.~Li, C.~Hu, et~al., ``Repeat-pass space-surface bistatic SAR tomography: accurate imaging and first experiment,'' \emph{Science China Information Sciences}, vol.~67, art.~192304, 2024. doi:10.1007/s11432-024-4089-2

\bibitem{ding2002} Y.~Ding, D.~C.~Munson~Jr., ``A fast back-projection algorithm for bistatic SAR imaging,'' in \emph{Proc. IEEE Int. Conf. on Image Processing (ICIP)}, vol.~II, pp.~449--452, 2002.

\bibitem{garza2010} G.~Garza, Z.~Qiao, ``Resolution analysis of bistatic SAR,'' in \emph{Proc. SPIE 8021, Radar Sensor Technology XV}, 80210F, 2010.

\bibitem{ge2019} N.~Ge, X.~X.~Zhu, ``Bistatic-like Differential SAR Tomography,'' \emph{IEEE Transactions on Geoscience and Remote Sensing}, vol.~57, no.~8, pp.~5883--5893, 2019.

\bibitem{grathwohl2023} A.~Grathwohl, B.~Meinecke, M.~Widmann, J.~Kanz, C.~Waldschmidt, ``UAV-Based Bistatic SAR-Imaging Using a Stationary Repeater,'' \emph{IEEE Journal of Microwaves}, vol.~3, no.~2, pp.~625--634, 2023. doi:10.1109/JMW.2023.3253667

\bibitem{he2025} J.~He, H.~Xie, H.~Liu, Z.~Wu, B.~Xu, N.~Zhu, Z.~Lu, P.~Qin, ``Multi-Baseline Bistatic SAR Three-Dimensional Imaging Method Based on Phase Error Calibration Combining PGA and EB-ISOA,'' \emph{Remote Sensing}, vol.~17, no.~3, art.~363, 2025. doi:10.3390/rs17030363

\bibitem{homer2005} J.~Homer, E.~Donskoi, B.~Mojarrabi, J.~Palmer, K.~Kubik, ``Three-dimensional bistatic synthetic aperture radar imaging system: spatial resolution analysis,'' \emph{IEE Proceedings -- Radar, Sonar and Navigation}, vol.~152, no.~6, pp.~391--394, 2005.

\bibitem{horne2002} A.~M.~Horne, G.~Yates, ``Bistatic synthetic aperture radar,'' in \emph{Proc. RADAR 2002}, Edinburgh, IEE Publication No.~490, pp.~6--10, 2002.

\bibitem{kanz2025} J.~Kanz, C.~Gesell, C.~Bonfert, D.~Werbunat, A.~Grathwohl, J.~Aguilar, M.~Vossiek, C.~Waldschmidt, ``UAV-Borne Digital Radar System for Coherent Multistatic SAR Imaging,'' arXiv:2507.20792, 2025.

\bibitem{liu2024} M.~Liu et~al., ``Unmanned Airborne Bistatic Interferometric Synthetic Aperture Radar Data Processing Method Using Bi-Directional Synchronization Chain Signals,'' \emph{Remote Sensing}, vol.~16, no.~5, art.~769, 2024. doi:10.3390/rs16050769

\bibitem{orfeo2019} D.~Orfeo, Y.~Zhang, D.~Burns, J.~S.~Miller, D.~R.~Huston, T.~Xia, ``Bistatic Antenna Configurations for Air-Launched Ground Penetrating Radar,'' \emph{Journal of Applied Remote Sensing}, vol.~13, no.~2, art.~027501, 2019. doi:10.1117/1.JRS.13.027501

\bibitem{ouchi2017} K.~Ouchi, C.-S.~Yang, ``Direct Evidence and Implications of Brewster's Angle Damping Observed by Synthetic Aperture Radar Images,'' in \emph{Proc. IGARSS 2017}, 2017.

\bibitem{ouchi2018} K.~Ouchi, C.-S.~Yang, ``Brewster Angle Damping Observed in the TerraSAR-X Synthetic Aperture Radar Images of Man-Made Targets,'' \emph{IEEE Geoscience and Remote Sensing Letters}, vol.~15, no.~4, pp.~532--536, 2018.

\bibitem{sun2022} H.~Sun, Z.~Sun, T.~Chen, Y.~Miao, J.~Wu, J.~Yang, ``An Efficient Backprojection Algorithm Based on Wavenumber-Domain Spectral Splicing for Monostatic and Bistatic SAR Configurations,'' \emph{Remote Sensing}, vol.~14, no.~8, art.~1885, 2022. doi:10.3390/rs14081885

\bibitem{thirion2018} L.~Thirion-Lefevre, R.~Guinvarc'h, ``The Double Brewster Angle Effect,'' \emph{Comptes Rendus Physique}, vol.~19, no.~1--2, pp.~43--53, 2018. doi:10.1016/j.crhy.2018.02.003

\bibitem{wang2020} Z.~Wang, M.~Liu, G.~Ai, P.~Wang, K.~Lv, ``Focusing of Bistatic SAR With Curved Trajectory Based on Extended Azimuth Nonlinear Chirp Scaling,'' \emph{IEEE Transactions on Geoscience and Remote Sensing}, vol.~58, no.~6, pp.~4160--4179, 2020.

\bibitem{wu2025} J.~Wu, J.~Yang, Z.~Li, W.~Pu, Z.~Sun, H.~An, Q.~Yang, ``Review of Bistatic Synthetic Aperture Radar Imaging Methods,'' \emph{Journal of Radars}, vol.~14, no.~5, pp.~1115--1141, 2025. doi:10.12000/JR25067

\bibitem{zhang2018} H.~Zhang, J.~Tang, R.~Wang, Y.~Deng, W.~Wang, N.~Li, ``An Accelerated Backprojection Algorithm for Monostatic and Bistatic SAR Processing,'' \emph{Remote Sensing}, vol.~10, no.~1, art.~140, 2018. doi:10.3390/rs10010140

\bibitem{zhao2025} T.~Zhao, J.~Wang, Z.~Cheng, Z.~Huang, J.~Song, ``A Subimage Autofocus Bistatic Ground Cartesian Back-Projection Algorithm for Passive Bistatic SAR Based on GEO Satellites,'' \emph{Remote Sensing}, vol.~17, no.~9, art.~1576, 2025. doi:10.3390/rs17091576

\end{thebibliography}

\end{document}
